\مسئله{با استفاده از مکان قطب‌ها و صفرها نسبت به دایره‌ی واحد، یک فیلتر آل‌پس طراحی کنید؛ به‌طوری‌که قطب و صفر سیگنال در بالای محور افقی قرار داشته باشند.}

\حل{
برای طراحی یک فیلتر آل‌پس دیجیتال با استفاده از مکان قطب‌ها و صفرها، ویژگی‌های زیر را مد نظر قرار می‌دهیم:
\begin{enumerate}
    \item \textbf{پاسخ فرکانسی ثابت:} اندازه پاسخ فرکانسی یک فیلتر آل‌پس برای تمام فرکانس‌ها برابر با یک (یا یک مقدار ثابت) است: $|H(e^{j\omega})| = 1$.
    \item \textbf{رابطه قطب و صفر:} برای تحقق این ویژگی، صفرها و قطب‌های فیلتر باید نسبت به دایره واحد هندسی معکوس و مزدوج یکدیگر باشند. یعنی اگر فیلتر دارای قطبی در نقطه $z_p = r e^{j\theta}$ باشد، باید صفری در نقطه زیر داشته باشد:
    $$z_z = \frac{1}{z_p^2} = \frac{1}{r} e^{j\theta}$$
    \item \textbf{پایداری سیستم:} برای پایدار بودن فیلتر زمان‌گسسته‌ و علی، تمام قطب‌ها باید درون دایره واحد قرار داشته باشند ($r < 1$). در نتیجه، صفرها خارج از دایره واحد ($1/r > 1$) خواهند بود.
    \item \textbf{موقعیت سیگنال (بالای محور افقی):} از آنجا که ذکر شده سیگنال بالای محور افقی قرار دارد، فاز قطب و صفر ($\theta$) باید بین $0$ تا $\pi$ (رادیان) انتخاب شود تا قطب و صفر در نیم‌صفحه بالایی دایره واحد قرار گیرند ($0 < \theta < \pi$).
\end{enumerate}

\paragraph{تابع تبدیل فیلتر:}
با فرض یک فیلتر آل‌پس مرتبه اول با قطب و صفر مختلط در نیم‌صفحه بالایی، تابع تبدیل به صورت زیر تعریف می‌شود:
$$H(z) = \frac{z^{-1} - z_z^{-1}}{1 - z_p z^{-1}}$$

با جایگذاری $z_p = r e^{j\theta}$ و صفر متناظر آن، تابع تبدیل نهایی به شکل زیر به دست می‌آید:
$$H(z) = e^{-j\theta} \frac{z^{-1} - r e^{-j\theta}}{1 - r e^{j\theta} z^{-1}}$$

اگر بخواهیم ضرایب فیلتر حقیقی باشند، باید زوج قطب و صفر مزدوج پیچیده را اضافه کنیم که در این صورت فیلتر مرتبه دوم خواهد شد:
$$H(z) = \frac{r^2 - 2r\cos(\theta)z^{-1} + z^{-2}}{1 - 2r\cos(\theta)z^{-1} + r^2 z^{-2}}$$
که در آن $0 < r < 1$ و $0 < \theta < \pi$ است تا قطب‌ها در داخل دایره واحد و در بالای محور افقی قرار بگیرند.
}